\clearpage\section*{阅读前：证明中反复使用的语言}
\addcontentsline{toc}{section}{阅读前：证明中反复使用的语言}
这里先把后文使用的符号和操作说清楚，避免把读不懂的符号误当成已经完成的一步证明。正式条目中的“来源页码”只用于查证出处，不要求你打开原课件才能继续读。

\paragraph{任取、选取与固定}
证明 $A\subseteq B$，先\emph{任取} $x\in A$，说明论证必须适用于每个可能的 $x$，然后由已知条件推出 $x\in B$。而“\emph{选取}一个 $y$”通常依赖存在性条件，例如集合非空、并集归属、确界逼近。每次选取都须说明为什么能选。“固定”意味着后面的变化中暂不改它；例如固定 $x_0,\eps$ 后选出的邻域可以依赖它们，但不能再依赖附近任意点 $x$。

\paragraph{等号与包含的核查}
$A=B$ 要两向包含；只有 $A\subseteq B$ 不能消去另一侧。对每个 $x$ 的双向逻辑等价 $x\in A\iff x\in B$ 也能一次证明相等。链式推导中每个 $\iff$ 都可反向，单向蕴含不能写成 $\iff$。集合补集用 $CA$ 或 $A^c$，始终相对同一个指定全集；集合名字上的撇号 $A'$ 一般只是另一个集合，除非题目另外定义，不能自动读作补集。

\paragraph{并、交与量词}
\[
x\in\bigcup_iA_i\iff\exists i:x\in A_i,
\qquad
x\in\bigcap_iA_i\iff\forall i:x\in A_i.
\]
否定全称变成存在反例，否定存在变成全部不满足。特别是“每个集合各有一个点”不等于“有同一个点属于全部集合”；这是对应逆像伪证明中多次出错的根源。对称差 $A\triangle B=(A\setminus B)\cup(B\setminus A)$ 只保留恰在一边的元素。

\paragraph{数列与子列}
$x_n\to x$ 意味任给 $\eps>0$，存在 $N$，使所有 $n\ge N$ 都满足 $\|x_n-x\|<\eps$。子列 $x_{n_k}$ 要求下标严格递增，因此 $n_k\to\infty$。原列已收敛时每条子列保留同一极限；反过来只有某条子列收敛，不保证原列收敛。有限次嵌套抽子列时，要把所有相关量一起限制到同一下标，不得独立抽后随意拼接。

\paragraph{如何分配误差}
三角不等式 $\|x-z\|\le\|x-y\|+\|y-z\|$ 将总误差拆成两段。若目标小于 $\eps$，可把每段控制在 $\eps/2$；三段可各用 $\eps/3$。对有限 $m$ 项，可以先各控 $\eps/m$ 再取有限个起始指标的最大值。这个最大值使\emph{全部}条件从同一项以后同时有效。

\paragraph{开球、闭包、内部和相对概念}
$B(x,r)=\{y:\|y-x\|<r\}$。开集要求每个成员有一个完全落在集合内的正半径球。$x\in\overline A$ 要求每个球都遇到 $A$；在度量空间等价于有 $a_n\in A$ 趋于 $x$。闭集包含所有集内收敛列的极限。内部 $A^\circ$ 收集可被小球完全包在集合内的点，边界 $\partial A=\overline A\setminus A^\circ$。若背景空间为 $S$，相对开、相对闭及补集都在 $S$ 中计算，不要求它们在整个欧氏空间有同一性质。

\paragraph{紧性在证明里怎样用}
紧性定义是每个开覆盖能取有限子覆盖。有限维中它等价于闭且有界，也等价于每个集内序列有收敛到集内的子列。使用时要明确：开覆盖的是哪个集？抽的序列在哪个紧集中？取到的极限为什么仍在集内？另一种等价形式是：紧集内的一族闭子集若任意有限交非空，则全部交非空。因为若全部交空，它们的相对开补集覆盖紧集，有限子覆盖就会使某个有限闭交空，矛盾。

\paragraph{集合值对应的三种检查}
$\Phi:S\rightrightarrows T$ 的值 $\Phi(x)$ 是集合；$y\in\Phi(x)$ 是输出点的归属。图为 $\{(x,y):y\in\Phi(x)\}$。闭图检查\emph{已有收敛选择}的极限不漏出；下半连续保证旧选择可以被附近的选择逼近；上半连续保证包含旧值集的每个开集最终也包含附近全部值集。仅闭值、闭图、紧值和上半连续是不同条件，不可混称。

\paragraph{引用与补充的边界}
本册只重写老师实际展示了过程的证明。课程仅陈述的深定理不会被偷偷添加成新的教师证明；若某一步引用它，会写明用途。遇到原证明错误，正文会指出第一次失效的地方，再给符合原路线的补正或反例，而不把无效公式当成合法一步。较短的正确证明不靠重复文字拉长，较长的证明也不以“同理”“显然”代替关键选点与极限论证。
